% Generated draft: requires mathematical and source review.
\documentclass[11pt]{article}
\usepackage[margin=25mm]{geometry}
\usepackage{fontspec}
\setmainfont{Latin Modern Roman}
\usepackage{amsmath,amssymb,mathtools,unicode-math}
\setmathfont{Latin Modern Math}
\usepackage{enumitem,needspace,xurl,hyperref}
\setlength{\parindent}{0pt}
\setlength{\parskip}{5pt}
\setlength{\emergencystretch}{3em}
\newcommand{\N}{\mathbb{N}}
\newcommand{\Z}{\mathbb{Z}}
\newcommand{\Q}{\mathbb{Q}}
\newcommand{\R}{\mathbb{R}}
\newcommand{\C}{\mathbb{C}}
\newcommand{\F}{\mathbb{F}}
\newcommand{\A}{\mathbb{A}}
\newcommand{\PP}{\mathbb{P}}
\newcommand{\E}{\mathbb{E}}
\newcommand{\eps}{\varepsilon}
\newcommand{\dd}{\,\mathrm{d}}
\begin{document}
\section*{UnsolvedMath ID: 30001134}
\textbf{Source problem code:} catalogue entry
\subsection{Definitions and mathematical statement}
For each integer \(d\geq 1\), consider the following precise candidate interpretation of the source question. Write
\[
S^d=\{x\in\R^{d+1}:\|x\|=1\}.
\]
Let \(K,\tau)\) consist of a finite simplicial complex \(K\) and a simplicial involution \(\tau\), with a homeomorphism \(h:|K|\to S^d\) satisfying
\[
h\circ|\tau|=-h.
\]
Thus the symmetry corresponds to the antipodal map on the sphere. For a vertex \(v\), define its degree and the maximum vertex degree by
\[
\deg_K(v)=\#\{w\in V(K)\setminus\{v\}:\{v,w\}\in K\},
\qquad
\Delta(K)=\max_{v\in V(K)}\deg_K(v).
\]
Under this interpretation, set
\[
f_{\mathrm{sph}}(d)=\min_{(K,\tau)}\Delta(K).
\]
For the polytope formulation, let \(\mathcal P_d\) be the class of full-dimensional convex simplicial polytopes \(P\subset\R^{d+1}\) with \(P=-P\); translating the center to the origin does not change vertex degrees. Define
\[
f_{\mathrm{poly}}(d)=\min_{P\in\mathcal P_d}\Delta(\partial P),
\]
where \(\partial P\) denotes the boundary simplicial complex. These classes are nonempty, as witnessed by the boundary of the \((d+1)\)-dimensional cross-polytope. Every such polytope boundary gives an antipodally symmetric sphere triangulation, so
\[
f_{\mathrm{sph}}(d)\leq f_{\mathrm{poly}}(d).
\]
\textbf{Requested determination and unresolved scope.} Determine the minimum possible maximum vertex degree as a function of \(d\). The supplied record calls this function \(f(d)\) and describes the sphere and polytope formulations as equivalent. It does not specify whether sphere triangulations are restricted to polytopal ones, nor does it supply a justification that the two minima agree. Consequently, the definitions above are explicit candidate formulations, and their equality is not asserted here as a verified fact. The intended meaning of central symmetry and the intended class of triangulations require confirmation against the source.

The record reports
\[
f(1)=2,\qquad f(2)=4,\qquad d+1<f(d)\leq 2d.
\]
These are supplied claims, not independently verified results in this draft. The strict lower bound contradicts the reported value at \(d=1\). Its intended dimension range is missing; no corrected range is silently imposed. Positive dimensions are used for the candidate formulations above because the record explicitly discusses \(d=1,2\); it gives no convention for \(d=0\).
\subsection{Short English statement}
Determine the smallest achievable maximum vertex degree among centrally symmetric triangulations of a sphere. The source also gives a simplicial polytope formulation, but the claimed equivalence needs clarification. It reports values in dimensions \(1\) and \(2\) and general bounds whose stated range conflicts with the dimension-\(1\) value. The supplied dated curation labels the problem open; this is a statement-only draft requiring human mathematical review.
\subsection{Sources}
\begin{itemize}
\item Supplied local UnsolvedMath registry record, permanent ID \(30001134\), source problem code \texttt{catalogue entry}, entitled \emph{Minimum Maximum Degree of Centrally Symmetric Sphere Triangulations}, in the Oberwolfach Reports Open Problems collection. The record gives proposed year \(2009\) and names no individual proposer.
\item Supplied report citation: \emph{Discrete Differential Geometry} (2009), Oberwolfach Reports, pp.~75--144, DOI \texttt{10.4171/owr/2009/02}. \url{https://doi.org/10.4171/owr/2009/02}. The supplied background describes Oberwolfach Reports as documenting mathematical workshops and their open-problem discussions; no problem-specific page is supplied.
\item Supplied literature source: Rote, \emph{Open Problems in Discrete Differential Geometry}. \url{https://page.mi.fu-berlin.de/rote/Kram/OWR-DDG09-problems.pdf}.
\item Supplied literature source labeled ``original OWR report'': \url{https://doi.org/10.4171/owr/2009/02}.
\item Literature assessment dated exactly \texttt{2026-08-21}: the supplied curation labels the problem open, says that the exact function remains undetermined in the checked literature, and reports \(d+1<f(d)\leq2d\). This is supplied dated curation derived from the OWR material, not an independently verified literature assessment. Its time-sensitive status and cited claims require independent verification.
\end{itemize}
\end{document}
